% --------------------------------------------------------------

\begin{frame}{Narzędzia -- gadżety}

\begin{columns}
\column{0.4\textwidth}
    Konstruujemy relację przy pomocy:
    \begin{itemize}
        \item relacji krawędzi \(\to\)
        \item koniunkcji \(\land\)
        \item \textcolor{cfblue}{równości wierzchołków \(=\)}
        \item kwantyfikatora \(\exists\)
    \end{itemize}

\column{0.6\textwidth}
\begin{figure}
    \centering
    \begin{tikzpicture}[
    line width=1pt,
]
\begin{scope}[shift={(-1,0)}]

    \node[circle, line width=1pt, draw=white, fill=black] (v1) at (0,2) {};
    \node[above left] at (v1) {\(v_1\)};
    
    \node[circle, line width=1pt, draw=white, fill=black] (v2) at (2,2) {};
    \node[above right] at (v2) {\(v_2\)};

    \node[circle, line width=1pt, draw=white, fill=black] (v3) at (2,0) {};
    \node[below right] at (v3) {\(v_3\)};

    \node[circle, line width=1pt, draw=white, fill=black] (v4) at (0,0) {};
    \node[below left] at (v4) {\(v_4\)};


    \draw[->] (v1) -- (v2);
    \draw[->] (v2) -- (v3);
    \draw[->] (v3) -- (v4);
    \draw[->] (v4) -- (v1);

\end{scope}

\begin{scope}[shift={(0,-2)}]
\node[] at (-2.4,1){\textcolor{cfblue}{Różnowartościowość:}};

\node[] at (0,0){\(I(v_1,\dots,v_4)\equiv \underbrace{v_1 \to v_2 \land v_2 \to v_3 \land v_3 \to v_4 \land v_4 \to v_1}_\text{wszystkie krawędzie z rdzenia} \)};
\end{scope}
\end{tikzpicture}
\end{figure}
\end{columns}

\end{frame}

% --------------------------------------------------------------

\begin{frame}{Narzędzia -- gadżety}

\begin{columns}
\column{0.4\textwidth}
    Konstruujemy relację przy pomocy:
    \begin{itemize}
        \item relacji krawędzi \(\to\)
        \item koniunkcji \(\land\)
        \item równości wierzchołków \(=\)
        \item \textcolor{cfblue}{kwantyfikatora \(\exists\)}
    \end{itemize}

\column{0.6\textwidth}
\begin{figure}
    \centering
    \begin{tikzpicture}[
    line width=1pt,
]

    \node[font=\Large] at (1.4, 3){\textcolor{cfblue}{\(R(u) \equiv \forall_{v \in V}\) \( u \to v\)}};
    \node[font=\Large] at (-1, 0.75){\(\exists_{v_1, \dots, v_4}\)};
    \node[font=\Large] at (2.75, 0.75){\( \land \bigwedge_{i=1}^4 u \to v_i\)};
        
    \node[circle, line width=1pt, draw=white, fill=black] (v1) at (0,1.25) {};
    \node[above left] at (v1) {\(v_1\)};
    
    \node[circle, line width=1pt, draw=white, fill=black] (v2) at (1,1.25) {};
    \node[above right] at (v2) {\(v_2\)};

    \node[circle, line width=1pt, draw=white, fill=black] (v3) at (1,0.25) {};
    \node[below right] at (v3) {\(v_3\)};

    \node[circle, line width=1pt, draw=white, fill=black] (v4) at (0,0.25) {};
    \node[below left] at (v4) {\(v_4\)};


    \draw[->] (v1) -- (v2);
    \draw[->] (v2) -- (v3);
    \draw[->] (v3) -- (v4);
    \draw[->] (v4) -- (v1);

\end{tikzpicture}
\end{figure}
\end{columns}

\end{frame}


% --------------------------------------------------------------

% \begin{frame}{Rdzeń -- dodatkowe utrudnienia}

% W równowartościowości wierzchołki nie mogą być na wszystkich pozycjach relacji.\\
% \vspace{1em}
% \textcolor{cfblue}{Kilka kopii rdzenia:} \\

% Różnowartościowość działa w obrębie kopii.

% \end{frame}

% --------------------------------------------------------------